% Week 1 section of the canonical synthesis exercise sheet.
% Included by Exercises.tex.

\section{Week 1 --- Option Pricing and Core Probability Tools}

\lecturesummary{Random variables, distributions, and characteristic functions;
sample statistics, the Law of Large Numbers, and the Central Limit Theorem;
joint distributions and information; conditional expectation; processes,
filtrations, and martingales; and a short differential-calculus review.
Exercises marked ``Extension'' may be left for additional practice.}

\subsection*{Part I --- Random Variables and Distributions}

\begin{exercise}[Level 1 --- A financial payoff as a random variable]
At maturity $T=1$, a stock price has the physical distribution
\[
\begin{array}{c|ccc}
s & 70 & 105 & 140\\ \hline
\Pb(S_T=s) & 0.25 & 0.50 & 0.25
\end{array}
\]
and $X=(S_T-100)^+$ is the payoff of a European call.
\begin{exparts}
  \item Identify the underlying, maturity, strike, and possible values of $X$.
  \item Determine the law and the cdf of $X$.
  \item Compute $\E[X]$, $\E[X^2]$, $\operatorname{Var}(X)$, and
  $\operatorname{SD}(X)$.
  \item If $r=3\%$ is continuously compounded, compute $e^{-rT}\E[X]$.
  Explain why this discounted physical expectation is not, by itself, a
  no-arbitrage price.
\end{exparts}
\end{exercise}

\begin{exercise}[Level 1 --- Cdf, density, and moments]
Let $U\sim\operatorname{Unif}[-1,3]$.
\begin{exparts}
  \item Write the cdf $F_U$ on the whole real line and recover the density $f_U$.
  \item Compute $\Pb(0<U\leq2)$ from the cdf and again from the density.
  \item Starting from the defining integrals, compute $\E[U]$, $\E[U^2]$,
  and $\operatorname{Var}(U)$.
  \item Set $V=2U+1$. Find its support, cdf, mean, and variance.
\end{exparts}
\end{exercise}

\begin{exercise}[Level 2 --- Empirical moments, LLN, and CLT]
The observed daily claim counts in a short sample are
\[
2,\ 5,\ 4,\ 3,\ 6.
\]
In the population, suppose daily claim counts $X_1,X_2,\ldots$ are i.i.d.
Poisson with parameter $4$.
\begin{exparts}
  \item Compute the sample mean, the empirical variance with denominator $n$,
  and the unbiased sample variance with denominator $n-1$.
  \item State the LLN for $\overline X_n$ and interpret it here.
  \item State the CLT for $\overline X_n$. What are the approximate mean and
  standard deviation of $\overline X_{100}$?
  \item Use the CLT to approximate
  $\Pb(\overline X_{100}>4.3)$. You may use $\Phi(1.5)\approx0.9332$.
  \item Explain in one sentence the different conclusions supplied by the LLN
  and the CLT.
\end{exparts}
\end{exercise}

\begin{exercise}[Level 2 --- Characteristic function of a Gaussian]
Let $Z\sim\mathcal N(0,1)$ and define
\[
\varphi_Z(u)=\E[e^{iuZ}],\qquad u\in\R.
\]
The following steps derive the characteristic function without using a
complex change of variables.
\begin{exparts}
  \item Using the standard Gaussian density, write $\varphi_Z(u)$ as an
  integral and verify that $\varphi_Z(0)=1$.
  \item Differentiate under the integral sign to express
  $\varphi_Z'(u)$. You may use
  \[
  \frac{\dd}{\dd z}e^{-z^2/2}=-z e^{-z^2/2}.
  \]
  \item Integrate by parts and show that
  \[
  \varphi_Z'(u)=-u\varphi_Z(u).
  \]
  Check explicitly that the boundary term vanishes.
  \item Solve this differential equation using the initial condition from
  part (a), and obtain $\varphi_Z(u)$.
  \item If $X=\mu+\sigma Z\sim\mathcal N(\mu,\sigma^2)$, deduce that
  \[
  \varphi_X(u)=\exp\!\left(i\mu u-\frac12\sigma^2u^2\right).
  \]
  \item Let $X_1$ and $X_2$ be independent Gaussian variables. Use
  characteristic functions to identify the distribution of $X_1+X_2$.
\end{exparts}
\end{exercise}

\subsection*{Part II --- Joint Distributions and Information}

\begin{exercise}[Level 1 --- Reading a joint law]
The joint law of $(X,Y)$ is
\[
\begin{array}{c|cc|c}
 & Y=L & Y=H & \Pb(X=x)\\ \hline
X=0 & 0.10 & 0.40 & {}\\
X=1 & 0.20 & 0.30 & {}\\ \hline
\Pb(Y=y) & {} & {} & 1
\end{array}
\]
\begin{exparts}
  \item Complete the marginal probabilities in the table.
  \item Are $X$ and $Y$ independent? Give a numerical justification.
  \item Compute $\Pb(X=1\mid Y=L)$ and $\Pb(X=1\mid Y=H)$.
  \item Interpret $X$ as a loss indicator and $Y$ as a rating signal. Which
  signal carries the larger conditional loss probability?
\end{exparts}
\end{exercise}

\begin{quote}\small
For a signal $Y$, $\sigma(Y)$ denotes exactly the information revealed by
observing $Y$. On a finite space, its events are the unions of the level sets
of $Y$.
\end{quote}

\begin{exercise}[Level 2 --- Information generated by a signal]
A fair die is rolled and $Y=\1_{\{4,5,6\}}$. Let $\mathcal G=\sigma(Y)$.
\begin{exparts}
  \item Give the partition of $\Omega=\{1,\ldots,6\}$ generated by $Y$ and
  list all events in $\mathcal G$.
  \item Decide whether each variable is $\mathcal G$-measurable:
  \[
  Z_1=Y,\qquad Z_2=(-1)^Y,\qquad
  Z_3=\1_{\{5,6\}},\qquad Z_4=2+3Y.
  \]
  Justify each answer using constancy on the atoms of the partition.
  \item Find $g:\{0,1\}\to\R$ such that $Z_2=g(Y)$.
  \item Explain what information is missing when one observes $Y$ but wants
  to determine $Z_3$.
\end{exparts}
\end{exercise}

\subsection*{Part III --- Conditional Expectation}

\begin{exercise}[Level 1 --- Conditional expectation on a finite partition]
A fair die is rolled. The available information reveals only the block
\[
L=\{1,2\},\qquad M=\{3,4\},\qquad H=\{5,6\},
\qquad \mathcal G=\sigma(L,M,H).
\]
Define $X=6\1_{\{\omega\geq4\}}$.
\begin{exparts}
  \item Is $X$ $\mathcal G$-measurable? Explain.
  \item Compute $\E[X]$ and the average of $X$ on each atom $L,M,H$.
  \item Write $\E[X\mid\mathcal G]$ using indicator functions.
  \item Verify the defining averaging property on $L$, $M$, and $H$.
  \item Verify $\E[\E[X\mid\mathcal G]]=\E[X]$.
\end{exparts}
\end{exercise}

\begin{exercise}[Level 2 --- Conditional density and conditional moments]
The pair $(X,Y)$ has joint density
\[
f_{X,Y}(x,y)=(x+y)\1_{\{0<x<1,\ 0<y<1\}}.
\]
\begin{exparts}
  \item Verify that the density integrates to one and compute both marginal densities.
  \item For $0<y<1$, derive $f_{X\mid Y}(x\mid y)$ and its support. Are
  $X$ and $Y$ independent?
  \item Compute $\E[X\mid Y=y]$, then write $\E[X\mid Y]$ as a random variable.
  \item For $f(x)=x^2$, compute $\E[f(X)\mid Y=y]$. Identify $g$ such that
  $\E[f(X)\mid Y]=g(Y)$ and explain why this is $\sigma(Y)$-measurable.
  \item Verify the tower property by comparing $\E[\E[X\mid Y]]$ with $\E[X]$.
\end{exparts}
\end{exercise}

\begin{exercise}[Extension --- Gaussian conditioning]
Let $(X,Y)$ be jointly Gaussian with
\[
\E[X]=\mu_X,\quad \E[Y]=\mu_Y,\quad
\operatorname{Var}(X)=\sigma_X^2,\quad
\operatorname{Var}(Y)=\sigma_Y^2>0,
\]
and $\operatorname{Cov}(X,Y)=\rho\sigma_X\sigma_Y$.
\begin{exparts}
  \item Set $a=\operatorname{Cov}(X,Y)/\operatorname{Var}(Y)$ and
  $R=X-\mu_X-a(Y-\mu_Y)$. Show that $\operatorname{Cov}(R,Y)=0$.
  \item Use joint Gaussianity to conclude that $R$ and $Y$ are independent.
  \item Deduce formulas for $\E[X\mid Y]$ and $\operatorname{Var}(X\mid Y)$.
  \item Interpret $|\rho|$ as hedge effectiveness and the conditional
  variance as residual risk.
\end{exparts}
\end{exercise}

\subsection*{Part IV --- Stochastic Processes and Filtrations}

\begin{exercise}[Level 2 --- Three coin tosses viewed through time]
Three fair coins are tossed. Write
$\omega=(\varepsilon_1,\varepsilon_2,\varepsilon_3)\in\{H,T\}^3$ and define
\[
\mathcal F_0=\{\varnothing,\Omega\},\qquad
\mathcal F_n=\sigma(\varepsilon_1,\ldots,\varepsilon_n),\quad n=1,2,3.
\]
Let $S_n$ be the number of heads observed by time $n$.
\begin{exparts}
  \item Describe the atoms of $\mathcal F_0,\ldots,\mathcal F_3$ and verify
  that these sigma-algebras form a filtration.
  \item Show that $(S_0,S_1,S_2,S_3)$ is adapted.
  \item Define $A_0=A_1=0$, $A_2=\1_{\{\varepsilon_3=H\}}$, and $A_3=A_2$.
  At which time does $(A_n)$ fail to be adapted? Interpret the failure.
  \item Let $\tau$ be the first head, capped at $3$. Determine $\tau$ on every
  outcome and verify $\{\tau\leq n\}\in\mathcal F_n$.
\end{exparts}
\end{exercise}

\subsection*{Part V --- Martingales and Fair Games}

\begin{exercise}[Level 2 --- A fair but asymmetric random walk]
Let $(\xi_k)_{k\geq1}$ be independent with
\[
\Pb(\xi_k=2)=\frac13,\qquad \Pb(\xi_k=-1)=\frac23,
\qquad S_n=\sum_{k=1}^n\xi_k,
\qquad \mathcal F_n=\sigma(\xi_1,\ldots,\xi_n).
\]
\begin{exparts}
  \item Compute $\E[\xi_k]$ and $\E[\xi_k^2]$. Why is this a fair game?
  \item Prove from the definition that $(S_n)$ is a martingale.
  \item Prove that $S_n^2-2n$ is a martingale.
  \item If instead $\Pb(\xi_k=2)=1/4$, find the predictable drift and
  construct a centered martingale from $S_n$.
  \item Explain the casino interpretation of the martingale property.
\end{exparts}
\end{exercise}

\begin{exercise}[Level 2 --- Replication and risk-neutral pricing]
A stock has $S_0=50$ and, after one period,
\[
S_1=65\quad\text{or}\quad S_1=45.
\]
The risk-free gross return is $R=1.05$. A call has strike $K=55$ and payoff
$H=(S_1-K)^+$.
\begin{exparts}
  \item Build the state-by-state payoff table.
  \item Find the risk-neutral up probability $q$ from
  $S_0=R^{-1}\E^q[S_1]$ and check that $q\in(0,1)$.
  \item Compute $V_0=R^{-1}\E^q[H]$.
  \item Find $\Delta$ shares and a maturity cash position $B_1$ such that
  $\Delta S_1+B_1=H$ in both states.
  \item Check that $\Delta S_0+B_1/R=V_0$ and explain why the physical up
  probability does not determine this price.
  \item Relate the pricing identity to the martingale property of discounted
  asset prices under the risk-neutral probability.
\end{exparts}
\end{exercise}

\subsection*{Part VI --- Mathematical Toolbox}

\begin{exercise}[Levels 1--2 --- Derivatives, gradients, and the chain rule]
\begin{exparts}
  \item Starting from the difference quotient, compute the derivative of
  $g(x)=x^2$.
  \item For $f(x,y)=x^2y+e^y$, compute all first and second partial
  derivatives, the gradient, and the Hessian.
  \item Let $S$ be differentiable and let $V=V(t,s)$ be smooth. If
  $h(t)=V(t,S(t))$, derive $h'(t)$ from the multivariable chain rule.
  \item Apply the previous result to $V(t,s)=e^{-rt}s^2$.
  \item The payoff $v(s)=(s-K)^+$ is continuous but not differentiable at one
  point. Identify it and compute the derivative everywhere else.
\end{exparts}
\end{exercise}
